\SourceLeaf{332}%
und hieraus, in Gemässheit von \(t=\tau(1+Frr)\),
\[
\begin{aligned}
t=\tau\{1+\frac{k^4}{12aa(1-ee)}\mathord{.}rr+\tfrac1{24}\tau\tau\\
{}+\frac1{24k^4}(5ee+(4ee-14e^4)ss+5e^4s^4)\beta\beta\}\quad .\ .\quad(9)
\end{aligned}
\]
\Article{31}
Für die sieben Bestandtheile von \(G\) ergeben sich folgende Werthe:
\[
\begin{aligned}
1)\quad&+\frac{kk\sin T^2}{12aacc}\mathord{.}ss\\
2)\quad&+\frac{kk\sin T^2}{12aa(1-ee)cc}\mathord{.}(1-2eess+ees^4)\\
3)\quad&-\frac{eekk\sin T^2}{4aa(1-ee)}\mathord{.}ss\\
4)\quad&-\frac{3e^4kk\cos T^2}{4aa(1-ee)^2}\mathord{.}ccss\\
5)\quad&-\frac{kk\sin T^2}{24aacc}\mathord{.}ss\\
6)\quad&+\frac{eekk\sin T^2}{4aa(1-ee)}\mathord{.}ss\\
7)\quad&-\frac{eekk\cos T^2}{8aa(1-ee)^2}\mathord{.}(1-(2+2ee)ss+3ees^4)
\end{aligned}
\]
Hier destruiren die Theile 3 und 6 einander; die übrigen vereinigen sich, indem man einerseits 1, 2 und 5, andererseits 4 und 7 zusammenfasst, in
\[
\begin{aligned}
&+\frac{kk\sin T^2}{24aa(1-ee)cc}\mathord{.}(2+(1-5ee)ss+2ees^4)\\
&-\frac{eekk\cos T^2}{8aa(1-ee)^2}\mathord{.}(1-(2-4ee)ss-3ees^4)
\end{aligned}
\]
Das erste Glied verwandelt sich, da \(2+(1-5ee)ss+2ees^4\) mit \(2cckk+3(1-ee)ss\) identisch ist, in
\[
\frac{k^4\sin T^2}{12aa(1-ee)}+\frac{kk\sin T^2}{8aacc}\mathord{.}ss
\]
Lösen wir hier \(\frac{k^4\sin T^2}{12aa(1-ee)}\) in \(\frac{k^4}{12aa(1-ee)}-\frac{kk\cos T^2}{12aa(1-ee)}\mathord{.}(1-eess)\) auf, so gibt die Vereinigung aller Theile
\[
\begin{aligned}
G={}&\frac{k^4}{12aa(1-ee)}+\frac{kk\sin T^2\operatorname{tang}B^2}{8aa}\\
&-\frac{kk\cos T^2}{24aa(1-ee)^2}\mathord{.}(2+ee-(8ee-14e^4)ss-9e^4s^4)
\end{aligned}
\]
und hieraus, in Gemässheit von \(b=\beta(1+Grr)\),
\[
\begin{aligned}
b=\beta\{1+\frac{k^4}{12aa(1-ee)}\mathord{.}rr+\tfrac18\tau\tau\\
{}-\frac1{24k^4}\mathord{.}(2+ee-(8ee-14e^4)ss-9e^4s^4)\beta\beta\}\quad .\ .\quad(10)
\end{aligned}
\]
